\chapter{Square roots, whitening, trace, determinant, and singular values}

\section{The symmetric square root of an SPD matrix}

If
\begin{equation}
\G=\mathbf Q\diag(g_1,\ldots,g_n)\mathbf Q^\top,
\qquad g_j>0,
\end{equation}
then the symmetric square root is
\begin{equation}
\G^{1/2}=\mathbf Q\diag(\sqrt{g_1},\ldots,\sqrt{g_n})\mathbf Q^\top.
\end{equation}

Multiplying it by itself gives
\begin{align}
\G^{1/2}\G^{1/2}
&=\mathbf Q\diag(\sqrt{g_j})\mathbf Q^\top
\mathbf Q\diag(\sqrt{g_j})\mathbf Q^\top\\
&=\mathbf Q\diag(g_j)\mathbf Q^\top\\
&=\G.
\end{align}

\inwords The square root applies half of the stretching represented by \(\G\). Applying that half-transformation twice reconstructs the full covariance transformation.

The inverse square root is
\begin{equation}
\G^{-1/2}=\mathbf Q\diag(g_1^{-1/2},\ldots,g_n^{-1/2})\mathbf Q^\top.
\end{equation}

\section{Whitening}

Define a whitened displacement
\begin{equation}
\mathbf x=\G^{-1/2}\mathbf d.
\label{eq:whiteningDef}
\end{equation}
Equivalently,
\begin{equation}
\mathbf d=\G^{1/2}\mathbf x.
\label{eq:unwhiteningDef}
\end{equation}

If \(\mathbf d\) has covariance \(\G\), then
\begin{align}
\Cov(\mathbf x)
&=\Cov(\G^{-1/2}\mathbf d)\\
&=\G^{-1/2}\G\G^{-1/2}\\
&=\Ii.
\end{align}

\inwords Whitening measures every displacement in units of the available genetic standard deviation. After whitening, the genetic covariance ellipsoid becomes a sphere.

Substituting \cref{eq:unwhiteningDef} into the fitness quadratic gives
\begin{align}
\mathbf d^\top\Hh\mathbf d
&=(\G^{1/2}\mathbf x)^\top\Hh(\G^{1/2}\mathbf x)\\
&=\mathbf x^\top\G^{1/2}\Hh\G^{1/2}\mathbf x.
\end{align}
This motivates the operator
\begin{equation}
\Sop=\G^{1/2}\Hh\G^{1/2}.
\label{eq:SdefinitionEarly}
\end{equation}

\begin{geombox}
A unit vector in whitened space represents one genetic standard deviation of displacement. Therefore \(\mathbf x^\top\Sop\mathbf x\) is the fitness curvature experienced by a genetically standardised displacement. The operator \(\Sop\) is the fitness landscape viewed through the amount and correlation structure of standing genetic variation.
\end{geombox}

\section{Trace as total curvature-weighted variance}

The trace is the sum of diagonal elements:
\begin{equation}
\tr(\mathbf A)=\sum_iA_{ii}.
\end{equation}
It is invariant under cyclic permutation:
\begin{equation}
\tr(\mathbf A\mathbf B)=\tr(\mathbf B\mathbf A).
\end{equation}

For a zero-mean random vector \(\boldsymbol\varepsilon\) with covariance \(\Pp\),
\begin{align}
\E[\boldsymbol\varepsilon^\top\Hh\boldsymbol\varepsilon]
&=\E\left[\tr(\boldsymbol\varepsilon^\top\Hh\boldsymbol\varepsilon)\right]\\
&=\E\left[\tr(\Hh\boldsymbol\varepsilon\boldsymbol\varepsilon^\top)\right]\\
&=\tr\left(\Hh\E[\boldsymbol\varepsilon\boldsymbol\varepsilon^\top]\right)\\
&=\tr(\Hh\Pp).
\label{eq:traceQuadraticExpectation}
\end{align}

\inwords The trace adds the variance in every direction after weighting each direction by fitness curvature. This identity is the mathematical bridge from hybrid covariance to hybrid fitness load.

\section{Determinant as volume scaling}

For a square transformation \(\mathbf A\), \(|\det\mathbf A|\) is the factor by which \(\mathbf A\) changes volume. For an SPD covariance matrix, \(\det\G\) is the generalized variance: it measures the squared volume scale of the covariance ellipsoid.

This volume interpretation later explains why
\begin{equation}
\det(\Vv_1^\top\Vv_2)^2
\end{equation}
measures squared overlap of two channel subspaces.

\section{Singular values}

Any rectangular matrix has a singular value decomposition
\begin{equation}
\mathbf A=\mathbf U\boldsymbol\Sigma\mathbf V^\top,
\end{equation}
where \(\mathbf U\) and \(\mathbf V\) have orthonormal columns and \(\boldsymbol\Sigma\) contains nonnegative singular values.

\usual Eigenanalysis asks for directions preserved by one square operator. Singular-value analysis asks how strongly one orthonormal set projects onto another. This is why principal angles between subspaces are defined using an SVD rather than an eigenanalysis.

\begin{summarybox}
The square root \(\G^{1/2}\) converts genetically standardised coordinates back to phenotype space. The inverse square root whitens genetic covariance. The trace converts a covariance into an expected quadratic loss. The determinant measures volume. Singular values measure mutual projection between subspaces. These four pieces of machinery carry nearly the entire theory.
\end{summarybox}

\part{Deriving additive genetic channels}
